\documentclass[11pt,letterpaper]{article}
\usepackage[margin=1in]{geometry}\usepackage{amsmath,amssymb,booktabs,array,microtype}\usepackage[T1]{fontenc}\usepackage{lmodern}
\usepackage[colorlinks=true,linkcolor=blue,citecolor=blue,urlcolor=blue]{hyperref}
\setlength{\parindent}{0pt}\setlength{\parskip}{5pt}\sloppy
\title{The Coldness Is Dynamical:\\ How Cold a Self-Critical MOND Sector Ends Up,\\ and Why Nothing Has To Be Tuned}
\author{Carl P.\ Zimmerman\\ \small Briar Creek Tech\\ \small AI-assisted research programme; not peer reviewed}
\date{12 September 2026}
\begin{document}\maketitle
\begin{abstract}
A cuscuton-clock MOND sector whose clock runs faster than proper time is gradient-unstable at zero field gradient, but its own MOND nonlinearity turns the sound speed positive at a finite gradient $Y_*$, whose marginal state has $c_s^2=0$ and isotropic stress \cite{paper19}. That left one thing unestablished: whether the back-reaction actually drives the gradient to $Y_*$ and holds it there closely enough for the Lyman-$\alpha$ forest. We settle it. Exactly two rates act on the gradient variance $Y=\langle|\nabla_\perp\chi|^2\rangle$: expansion dilutes each mode's contribution as $a^{-2}$, and where $c_s^2<0$ a mode of physical wavenumber $k$ pumps it at $2|c_s|k$. Growth switches off the instant $Y$ passes $Y_*$, so the two balance at $|c_s(Y_{\rm eq})|=H/k$, giving a residual $c_s^2=(H/k)^2$. Integrating the back-reaction from initial gradients spanning eight orders of magnitude around $Y_*$, every trajectory lands on the same balance point to six digits from both sides; the residual is $1/\kappa^2$ with $\kappa=k/H$ to six digits across four decades and at every gradient-unstable epoch of the candidate's own history; and with a whole spectrum evolving together the shortest wavelength present ends up carrying essentially all the gradient variance. The Lyman-$\alpha$ bound $c_s^2\le10^{-9}$ is therefore met once the instability acts on comoving scales below about $0.8$~Mpc at $z=3$. The apparent requirement that the gradient track $Y_*$ to one part in $10^6$ turns out to be the same condition in different units: the attractor sits at whatever precision the available wavelength dictates, and nothing is tuned. The cost is that the residual is set by the ultraviolet reach of the instability, so the forest gate becomes a lower bound on that reach rather than a tuning of coefficients.
\end{abstract}

\section{What was owed}
The action studied here is the frozen cuscuton-clock MOND candidate of \cite{paper13}, whose sub-horizon clock-scalar sound speed at zero field gradient is $c_s^2=(1-s_0)m_{\rm rel}/(2-m_{\rm rel})$ \cite{paper17}, negative whenever the clock rate $s_0=\dot\tau$ exceeds unity. Carrying a background gradient through the same reduction \cite{paper19} gives
\begin{equation}
c_s^2(Y)=\frac{2P_X(1-D)-2s_0\mathcal{W}}{B(1-D)},\qquad D=\frac{2Q^2\mathcal{W}}{W(Y)},\qquad \mathcal{W}=W_Y(Y)\ \text{or}\ W_Y+2YW_{YY},
\end{equation}
which rises monotonically in the projected gradient invariant $Y=|\nabla_\perp\chi|^2$ and crosses zero at a unique $Y_*$. The same surface is an exact off-shell stress degeneracy, $\det(T^{\rm mixed}_{tx}-L\mathbb{I})=-s_0b^2W\mathcal{N}$ with $\mathcal{N}$ the numerator above, so there the sector's stress takes perfect-fluid form. Isotropic stress with non-propagating perturbations is what a clustering cold component is, and a Lean-certified necessity result in the same programme \cite{repo} says any completion must contain one.

That argument is about signs. Below $Y_*$ the sector is unstable so gradients grow; above it expansion dilutes them. It does not say where the gradient ends up, nor how closely it holds. The Lyman-$\alpha$ forest needs $c_s^2\le10^{-9}$, and the sound speed rises steeply out of the marginal point, so the question is sharp.

\section{Two rates}
The gradient variance is carried by a spectrum of modes, and exactly two rates act on it. A comoving mode's physical gradient redshifts, so its contribution to $Y$ falls as $a^{-2}$. Where $c_s^2<0$, a mode of physical wavenumber $k$ grows at $|c_s|k$ in the WKB limit, so its contribution grows at twice that. In e-folds $N=\ln a$, with $\kappa\equiv k/H$,
\begin{equation}
\frac{d\ln Y}{dN}=-2+2\kappa\,\big|c_s(Y)\big|_{-},
\label{flow}
\end{equation}
where the second term is present only while $c_s^2<0$. The growth term switches off the instant $Y$ passes $Y_*$, and the flow is monotone in $Y$ on either side, so (\ref{flow}) has a single fixed point and it is an attractor:
\begin{equation}
\boxed{\ \big|c_s(Y_{\rm eq})\big|=\frac{1}{\kappa}=\frac{H}{k},\qquad c_s^2(Y_{\rm eq})=\Big(\frac{H}{k}\Big)^{2}.\ }
\label{fp}
\end{equation}
The residual sound speed is set by the balance between instability growth and Hubble dilution. No coefficient of the action appears in it.

\section{Integration}
Equation (\ref{flow}) is integrated on the candidate's own frozen coefficient functions and its archived branch, imported read-only, at the gradient-unstable epochs where $s_0>1$.

\emph{The fixed point is an attractor.} At $\kappa=10^3$ and epoch $a=0.56$ ($s_0=1.173$, $Y_*/\ell=0.0174$), starting gradients of $10^{-6}$, $10^{-2}$, $0.5$ and $50$ times $Y_*$ all land on the same balance point, $Y_{\rm eq}/Y_*=0.999725$, agreeing to six digits and approached from both sides.

\emph{The residual obeys (\ref{fp}).} Solving $|c_s(Y_{\rm eq})|=1/\kappa$ directly:
\begin{table}[h]\centering\caption{Residual sound speed and distance from the critical gradient, epoch $a=0.56$.}\label{tab}
\begin{tabular}{cccc}\toprule
$\kappa=k/H$ & $|c_s^2|$ at the balance & $1/\kappa^2$ & $|\Delta Y/Y_*|$\\\midrule
$10^2$ & $1.000\times10^{-4}$ & $10^{-4}$ & $2.74\times10^{-2}$\\
$10^3$ & $1.000\times10^{-6}$ & $10^{-6}$ & $2.75\times10^{-4}$\\
$10^4$ & $1.000\times10^{-8}$ & $10^{-8}$ & $2.75\times10^{-6}$\\
$10^5$ & $1.000\times10^{-10}$ & $10^{-10}$ & $2.75\times10^{-8}$\\
$10^6$ & $1.000\times10^{-12}$ & $10^{-12}$ & $2.75\times10^{-10}$\\\bottomrule
\end{tabular}\end{table}
$c_s^2\kappa^2=1.000000$ to six digits across four decades, and $1.0000$ at all five gradient-unstable epochs of the branch.

\emph{The shortest mode wins.} Evolving five wavenumbers spanning two decades together under the common $c_s^2(Y)$, the shortest wavelength ends up carrying essentially all the gradient variance and the resulting sound speed matches the single-mode answer for that wavenumber. The residual can legitimately be read off the shortest scale available.

\section{The forest gate, converted}
Setting $c_s^2=10^{-9}$ in (\ref{fp}) gives $\kappa=3.2\times10^4$. At $z=3$, with $H(z)/c=1.0\times10^{-3}\,{\rm Mpc}^{-1}$, that is a comoving wavelength of $0.78$~Mpc. The Lyman-$\alpha$ bound is therefore met once the instability acts on sub-megaparsec scales --- coarse compared with every other scale in this effective theory.

This also explains what the earlier ``precision requirement'' was. Linearising $c_s^2$ about $Y_*$ turns $c_s^2\le10^{-9}$ into a demand that the gradient track $Y_*$ to a fractional precision of $2.7\times10^{-7}$. Table~\ref{tab} shows the attractor parks the gradient at exactly $2.75\times10^{-7}$ when $\kappa=3.2\times10^4$: the two statements agree to better than one part in $10^3$ because they are the same condition written in different units. The attractor automatically sits at whatever precision the available wavenumber dictates. Nothing is tuned, and in particular the logarithm margin of the closure need not be tuned to $10^{-8}$, which is how this gate was previously met.

\section{What it costs, and what is not settled}
The residual (\ref{fp}) depends on the largest wavenumber the instability reaches. The coldness of the sector is therefore ultraviolet-sensitive, and the forest gate has become a lower bound on the instability's reach rather than a tuning of coefficients. That is a better place to be, but it is not free: a theory whose instability were cut off above $0.8$~Mpc at $z=3$ would fail.

The treatment is a mean-field (Hartree) back-reaction in the WKB limit, not a lattice simulation. Mode coupling, phase decoherence and the nonlinear saturation of individual modes are not modelled, and the growth rate $|c_s|k$ --- exact in WKB --- is precisely the assumption that makes the answer ultraviolet-sensitive. Only the forest gate is evaluated; the acoustic peaks, lensing, the post-Newtonian limit and the amount of the sector are untouched here. And the candidate's own coefficient history still fails the forest on its early branch, where the clock runs slower than proper time, no instability exists, and the criticality that produces the coldness never operates. The statement that survives is conditional and sharp: a history whose clock runs faster than proper time at the epochs of interest is driven by its own MOND sector to a state with isotropic stress and a sound speed of order $(H/k_{\max})^2$, and that is a condition on the coefficient functions which can be imposed a priori rather than fitted.

\section{Reproducibility}
\texttt{fable\_independent\_2026/L194\_tracking\_dynamics.py} with its committed output and \texttt{L194\_results.json} (commit \texttt{9212f4498}); the preceding steps in \texttt{L192\_gradient\_criticality.py} and \texttt{L193\_stress\_degeneracy.py}; the Lean file \texttt{fable\_independent\_2026/lean\_2026/Mondlean.lean} (111 theorems, no \texttt{sorry}). Coefficient functions and branch imported read-only from the collaborating programme named in \cite{paper13}.

\begin{thebibliography}{9}
\bibitem{repo} The research repository, \url{https://github.com/carlzimmerman/zimmerman-formula} (2026).
\bibitem{paper13} C.~P.~Zimmerman, \emph{A Primordial-Clock Relativistic MOND Candidate: Status Report}, doi:10.5281/zenodo.22699828 (2026).
\bibitem{paper17} C.~P.~Zimmerman, \emph{The Clock Stability Theorem}, doi:10.5281/zenodo.22717950 (2026).
\bibitem{paper19} C.~P.~Zimmerman, \emph{Gradient-Driven Criticality}, doi:10.5281/zenodo.22727111 (2026).
\bibitem{cuscuton} N.~Afshordi, D.~J.~H.~Chung, G.~Geshnizjani, Phys.\ Rev.\ D \textbf{75}, 083513 (2007).
\end{thebibliography}
\end{document}
